% SEGMENT OLP-0005-S01
% Part: sets-functions-relations
% Chapter: sets
% Section: basics

\documentclass[../../../include/open-logic-section]{subfiles}

\begin{document}

\olfileid{sfr}{set}{bas}
\olsection{உறுப்புசார் சமத்துவம்}

பொருள்களின் ஒரு தொகுப்பை ஒரே பொருளாகக் கருதும்போது அது ஒரு \emph{கணம்}
எனப்படும். அக்கணத்தை உருவாக்கும் பொருள்கள் அதன் \emph{உறுப்புகள்} அல்லது
\emph{அங்கங்கள்} எனப்படும். $x$ என்பது $A$ என்னும் கணத்தின் !!a{element}
எனில் $x \in A$ என எழுதுகிறோம்; இல்லையெனில் $x \notin A$ என எழுதுகிறோம்.
!!{element}s எதுவும் இல்லாத கணம் \emph{வெற்றுக்} கணம் எனப்படும்;
அதனை ``$\emptyset$'' எனக் குறிக்கிறோம்.

% SEGMENT OLP-0005-S02
\begin{explain}
ஒரு கணத்தை எவ்வாறு \emph{குறிப்பிடுகிறோம்}, அதன் !!{element}s எந்த
\emph{வரிசையில்} அமைகின்றன, அல்லது அதன் !!{element}s ஒவ்வொன்றையும்
\emph{எத்தனை முறை} எண்ணுகிறோம் என்பவை பொருட்டல்ல. அதன் !!{element}s
எவை என்பது மட்டுமே பொருட்டாகும். இதைப் பின்வரும் கொள்கையாக வகுக்கிறோம்.
\end{explain}

% SEGMENT OLP-0005-S03
\begin{defn}[உறுப்புசார் சமத்துவம்]
  $A$, $B$ ஆகியவை கணங்கள் எனில், $A$ இன் ஒவ்வொரு !!{element}[um]
  $B$ இன் !!a{element}[pred] இருப்பதுடன், மறுதிசையிலும் இதே நிபந்தனை நிறைவேறவும்
  செய்யும்போது, அப்போதும் அப்போது மட்டுமே $A = B$.
\end{defn}

% SEGMENT OLP-0005-S04
உறுப்புசார் சமத்துவம் சில குறியீடுகளைப் பயன்படுத்த வழிவகுக்கிறது.
பொதுவாக, $a_{1}$, \dots, $a_{n}$ ஆகிய பொருள்கள் தரப்பட்டால்,
$\{a_{1}, \dots, a_{n}\}$ என்பது $a_1, \ldots, a_n$ ஆகியவற்றைத் தனது
!!{element}s[pred] கொண்ட \emph{அந்தக்} கணமாகும். இங்கு ``\emph{அந்தக்}''
என்பதை வலியுறுத்துகிறோம்; ஏனெனில் இத்தகைய கணம் \emph{ஒன்றே ஒன்று}
மட்டுமே இருக்க முடியும் என உறுப்புசார் சமத்துவம் கூறுகிறது.
உண்மையில், உறுப்புசார் சமத்துவம் பின்வருவதையும் நியாயப்படுத்துகிறது:
  \[
    \{a, a, b\} = \{a, b\} = \{b,a\}.
  \]
கணங்களைக் கருதும்போது, அவற்றின் !!{element}s அமையும் வரிசையையோ,
அவை எத்தனை முறை குறிப்பிடப்படுகின்றன என்பதையோ பொருட்படுத்துவதில்லை
என்பதை இது தெளிவாக்குகிறது.

% SEGMENT OLP-0005-S05
\begin{tagblock}{novice}
\begin{ex}
பல பொருள்கள் உங்களிடம் இருக்கும்போதெல்லாம், அவற்றை ஒன்றுதிரட்டி ஒரு
கணமாக அமைக்கலாம். எடுத்துக்காட்டாக, ரிச்சர்டின் உடன்பிறந்தோரை உறுப்புகளாகக்
கொண்ட கணத்தில் ஒருவர் உள்ளார்; அதனை $S=\{\textrm{ரூத்}\}$ என எழுதலாம்.
$4$ ஐவிடச் சிறிய நேர்ம முழுக்களின் கணம் $\{1, 2, 3\}$ ஆகும்;
அதனை $\{3, 2, 1\}$ என்றோ $\{1, 2, 1, 2, 3\}$ என்றோகூட எழுதலாம்.
உறுப்புசார் சமத்துவத்தின்படி இவை அனைத்தும் ஒரே கணமே.
ஏனெனில் $\{1, 2, 3\}$ இன் ஒவ்வொரு !!{element}[um]
$\{3, 2, 1\}$ இன் (மேலும் $\{1, 2, 1, 2, 3\}$ இன்) !!a{element}
ஆகும்; மறுதிசையிலும் இதுவே உண்மை.
\end{ex}
\end{tagblock}

% SEGMENT OLP-0005-S06
பல சமயங்களில், ஒரு கணத்தின் !!{element}s அனைத்துக்கும் பொதுவான ஒரு
பண்பைக் கொண்டு அக்கணத்தைக் குறிப்பிடுவோம். அதற்குப் பின்வரும் சுருக்கக்
குறியீட்டைப் பயன்படுத்துவோம்: $\Setabs{x}{\phi(x)}$.
இங்கு $\phi(x)$ என்பது, அக்கணத்தின் !!{element}s மத்தியில் $x$ இடம்பெற
வேண்டுமெனில் அது கொண்டிருக்க வேண்டிய பண்பைக் குறிக்கிறது.

% SEGMENT OLP-0005-S07
\begin{tagblock}{novice}
\begin{ex}
நமது எடுத்துக்காட்டில், $S$ ஐப் பின்வருமாறும் குறிப்பிட்டிருக்கலாம்:
\[
S = \Setabs{x}{x \text{ ரிச்சர்டின் உடன்பிறந்தவர்}}.
\]
\end{ex}
\end{tagblock}

% SEGMENT OLP-0005-S08
\begin{tagblock}{math}
\begin{ex}
ஒரு எண் அதன் தகு வகுத்திகளின் கூட்டுத்தொகைக்குச் சமமாக இருக்கும்போது,
அப்போதும் அப்போது மட்டுமே அது \emph{நிறைவெண்} எனப்படும்.
இங்கு தகு வகுத்திகள் என்பவை, அவ்வெண்ணை மீதியின்றி வகுக்கும்,
ஆனால் அவ்வெண்ணுக்குச் சமமாக இல்லாத எண்களாகும்.
எடுத்துக்காட்டாக, $6$ ஒரு நிறைவெண்; ஏனெனில் அதன் தகு வகுத்திகள்
$1$, $2$, $3$ ஆகியவை; மேலும் $6 = 1 + 2 + 3$.
உண்மையில், $10$ ஐவிடச் சிறிய நேர்ம முழுக்களில் $6$ மட்டுமே
நிறைவெண்ணாகும். ஆகவே, உறுப்புசார் சமத்துவத்தைப் பயன்படுத்தி,
பின்வருமாறு கூறலாம்:
\[
	\{6\} = \Setabs{x}{x\text{ நிறைவெண் மற்றும் }0 \leq x \leq 10}
\]
வலப்புறக் குறியீட்டை ``$x$ நிறைவெண்ணாகவும் $0 \leq x \leq 10$
என்ற நிபந்தனையை நிறைவேற்றுவதாகவும் உள்ள $x$ களின் கணம்'' என வாசிக்கிறோம்.
கணங்களைக் கருதும்போது, அவற்றை எவ்வாறு குறிப்பிடுகிறோம் என்பது
பொருட்டல்ல என்பதை இங்குள்ள சமத்துவம் உறுதிப்படுத்துகிறது.
மேலும் பொதுவாக, $\phi(x)$ என்ற நிபந்தனையை நிறைவேற்றும் $x$ களின்
கணம் எப்போதும் ஒன்றே ஒன்றாக இருக்கும் என்பதை உறுப்புசார் சமத்துவம்
உறுதிப்படுத்துகிறது. எனவே, $\Setabs{x}{\phi(x)}$ என்பதை
$\phi(x)$ என்ற நிபந்தனையை நிறைவேற்றும் $x$ களின் \emph{அந்தக்} கணம்
என அழைப்பதை உறுப்புசார் சமத்துவம் நியாயப்படுத்துகிறது.
\end{ex}
\end{tagblock}

% SEGMENT OLP-0005-S09
கணங்கள் ஒன்றே என்பதை நிறுவுவதற்கு உறுப்புசார் சமத்துவம் ஒரு வழியைத் தருகிறது:
$A = B$ என நிறுவ, $x \in A$ எனில் $x \in B$ என்பதையும்,
$y \in B$ எனில் $y \in A$ என்பதையும் நிறுவ வேண்டும்.

% SEGMENT OLP-0005-S10
\begin{prob}
வெற்றுக் கணம் அதிகபட்சம் ஒன்றே இருக்க முடியும் என நிறுவுக;
அதாவது, $A$, $B$ ஆகியவை !!{element}s இல்லாத கணங்கள் எனில்,
$A = B$ எனக் காட்டுக.
\end{prob}

\end{document}
